\scanpage{14}
\uline{Proposition} Soit $G$ \unsure{gpe|groupe} algébrique affine sur corps alg. clos $k$. Alors $\operatorname{Pic}(G)$ est fini. Si $G_{\mathrm{réd}}$ est \deleted{lisse et} \uline{résoluble}, \deleted{$\operatorname{Pic}(G)$ \unsure{est}} $\operatorname{Pic}(G)$ est nul. Si $G$ est lisse, \deleted{\illegible{1}}\inserted{connexe}, \deleted{$G$ est un $\operatorname{Pic}(G)$ est nul \unsure{ssi}} 
\[
\operatorname{Pic}(G)\simeq\operatorname{Pic}(G/B)/\operatorname{Im}D(T)
\]

\diagram{\begin{tikzcd}[ampersand replacement=\&,column sep=small,row sep=tiny]
R \& \subset \& B \& \subset \& G\arrow[r] \& G/B\\[-2pt]
R\cap T \& \& \overset{\cup}{T} \& \& \&
\end{tikzcd}}

[$B$ Borel, $T\subset B$ tore maximal,
\[
\underset{\substack{\scriptscriptstyle\parallel\\\operatorname{Hom}(B,G_m)}}{D(T)}
 \longrightarrow \operatorname{Pic}(G/B) \quad\text{hom caractéristique}\,]
\]

\diagram{\begin{tikzcd}[ampersand replacement=\&,row sep=small,column sep=small]
\& G/R \arrow[r] \& G/B\\
\deleted{T}\arrow[r] \& \overset{\cup}{T/R} \&\\
\& G\arrow[r,dash] \& G/B\arrow[uu]\\
\& \overset{\cup}{T} \&
\end{tikzcd}}

\diagram{\begin{tikzcd}[ampersand replacement=\&,row sep=small]
\operatorname{Pic}(G/B) \arrow[d,equal] \& \deleted{\unsure{\operatorname{Pic}}}D(T/R) \arrow[l,hook'] \arrow[d]\\
\operatorname{Pic}(G/B) \& \deleted{\unsure{\operatorname{Pic}}}D(T)\arrow[l]
\end{tikzcd}}

\diagram{\begin{tikzcd}[ampersand replacement=\&,row sep=tiny]
G/R=\operatorname{GP}(n)\arrow[r] \& \unsure{\mathcal D_n|\mathcal D(n)|\mathrm{Drap}|\text{Dém.}}\\
\deleted{\unsure{\operatorname{Gl}|\operatorname{G}}} \&\\
G=\operatorname{Gl}(n)\arrow[r] \& \mathcal D_n\\
\overset{\cup}{T(n)} \&
\end{tikzcd}}

Si $G$ est \deleted{\unsure{réductive|réductive et}} \inserted{semi-simple}, $\operatorname{Pic}(G)=0$ sss $G$ est simplement connexe.

\uline{Proposition} $G$ \deleted{\unsure{réductive sur|réductif un}} groupe alg affine sur corps $k$. Alors $\exists$ entier $n>0$ [égal au degré d'une extension convenable $k'$ de $k$] tel que $\unsure{n\cdot|n}\operatorname{Pic}(G)=0$. \deleted{Si $G_{\mathrm{réd}}$ \deleted{est} est \deleted{résoluble}\inserted{unipotent}, \unsure{alors} alors on peut prendre pour $n$ une puissance de la \unsure{caract.}}

On \unsure{a} les précisions suivantes
% archive sheet number: 13
